% -*- coding: utf-8 -*-
% Consolidated 13 standalone reviewed Gaokao calculus solution LaTeX sources.
% All original standalone document bodies are preserved, except each standalone \maketitle
% is replaced by a uniform \section* heading for compendium compatibility.
% Original source SHA IDs are recorded. Original individual files remain recoverable
% in Git history and snapshot-before-tex-merge-2026-10-27.
\documentclass[UTF8,a4paper,11pt]{ctexart}
\usepackage[margin=25mm,headheight=15pt]{geometry}
\usepackage{amsmath,amssymb,amsthm,fancyhdr,hyperref,xcolor}
\hypersetup{colorlinks=true,urlcolor=blue!55!black,linkcolor=black}
\setlength{\parindent}{2em}
\setlength{\parskip}{0.25em}
\linespread{1.13}
\pagestyle{fancy}
\fancyhf{}
\fancyhead[L]{\small 三千弱水 · 高考数学 · 导数独立解答}
\fancyhead[R]{\small 2010—2025 历年真题合订}
\fancyfoot[C]{\thepage}
\renewcommand{\headrulewidth}{0.3pt}
\newcommand{\qtitle}[1]{\par\bigskip\noindent{\large\bfseries #1}\par\medskip}
\newcommand{\sourceurl}[1]{\par\noindent{\small\color{black!65}来源：\href{#1}{固定版本 LaTeX 源码（点击查看）}}\par}
\begin{document}
\begin{center}
{\huge\bfseries 高考数学 · 导数真题独立证明合订本}\par\medskip
{\large 2010—2025 年 · 13 份来源证明（含 2024 年两题）}\par\bigskip
{\small 题目来源见各节；数学复核状态不等于 PDF 编译质检状态。}
\end{center}
\tableofcontents
\clearpage
% SOURCE_ENTRY path=derivative/reviewed/gaokao-2010-anhui-liberal-q20-solution.tex sha=2bf5f9578af0de48e37b2ea9852affc1cb8014b3
% ORIGINAL_PREAMBLE \documentclass[UTF8,a4paper,11pt]{ctexart}
% ORIGINAL_PREAMBLE \usepackage{amsmath,amssymb}
% ORIGINAL_PREAMBLE \usepackage[margin=2.6cm]{geometry}
% ORIGINAL_PREAMBLE \usepackage[hidelinks]{hyperref}
% ORIGINAL_PREAMBLE \title{2010年安徽高考文科第20题：导数与极值}
% ORIGINAL_PREAMBLE \author{三千弱水数学题库}
% ORIGINAL_PREAMBLE \date{2026年10月24日}
\clearpage
\phantomsection
\addcontentsline{toc}{section}{2010年安徽高考文科第20题：导数与极值}
\section*{2010年安徽高考文科第20题：导数与极值}
\noindent\textbf{出处：}\url{https://github.com/DxAThing/Gaokao-Math-Problems-Compilation/blob/f1ae169b58a75b3db6e35c670b3f0d5c78718188/content/2010/anhui_liberal.tex#L426}

设 $f(x)=\sin x-\cos x+x+1$，$0<x<2\pi$。

\paragraph{驻点}
\[
f'(x)=\cos x+\sin x+1
     =1+\sqrt{2}\sin\left(x+\frac{\pi}{4}\right).
\]
在 $(0,2\pi)$ 中解 $f'(x)=0$，得且仅得
$x=\pi$、$x=3\pi/2$。

\paragraph{单调性}
在 $(0,\pi)$、$(\pi,3\pi/2)$、$(3\pi/2,2\pi)$
三个区间中，导函数符号依次为正、负、正。
因此函数在第一、第三区间严格递增，
在第二区间严格递减。

\paragraph{极值}
函数在 $x=\pi$ 处取得严格极大值
\[
f(\pi)=\pi+2,
\]
在 $x=3\pi/2$ 处取得严格极小值
\[
f(3\pi/2)=\frac{3\pi}{2}.
\]
定义域端点不参与极值判断。

\bigskip
\noindent\textbf{质检说明：}数学推导独立验证通过；
本源码尚未进行实际 XeLaTeX 编译。
% END_SOURCE_ENTRY path=derivative/reviewed/gaokao-2010-anhui-liberal-q20-solution.tex

% SOURCE_ENTRY path=derivative/reviewed/gaokao-2010-beijing-science-q18-solution.tex sha=7e7d20a0489c1ebf2175c7faf5ed7e6aa71f63ae
% ORIGINAL_PREAMBLE \documentclass[UTF8,a4paper,11pt]{ctexart}
% ORIGINAL_PREAMBLE \usepackage{amsmath,amssymb}
% ORIGINAL_PREAMBLE \usepackage[margin=2.5cm]{geometry}
% ORIGINAL_PREAMBLE \usepackage[hidelinks]{hyperref}
% ORIGINAL_PREAMBLE \title{2010 年北京高考理科第 18 题\\含参对数函数的单调性}
% ORIGINAL_PREAMBLE \date{2026 年 10 月 24 日}
\clearpage
\phantomsection
\addcontentsline{toc}{section}{2010 年北京高考理科第 18 题 · 含参对数函数的单调性}
\section*{2010 年北京高考理科第 18 题\\含参对数函数的单调性}
\noindent\textbf{原题来源：}\url{https://github.com/DxAThing/Gaokao-Math-Problems-Compilation/blob/f1ae169b58a75b3db6e35c670b3f0d5c78718188/content/2010/beijing_science.tex#L352}
\par
函数 $f_k(x)=\ln(1+x)-x+\frac{k}{2}x^2$，其中 $k\ge0$，
定义域为 $(-1,+\infty)$。

\paragraph{切线}
\[
 f_k'(x)=\frac1{1+x}-1+kx
        =\frac{x(kx+k-1)}{1+x}.
\]
当 $k=2$ 时，$f_2(1)=\ln2$，$f_2'(1)=\frac32$。
在该点的切线方程为
\[
 \boxed{y-\ln2=\frac32(x-1)}.
\]

\paragraph{完整参数分类}
导数的分母 $1+x$ 恒为正。令 $r=(1-k)/k$（$k>0$）。
\begin{enumerate}
\item 当 $k=0$ 时，$f_k'=-x/(1+x)$。函数在
$(-1,0)$ 严格递增，在 $(0,+\infty)$ 严格递减。
\item 当 $0<k<1$ 时，$r>0$，导数在
$(-1,0),(0,r),(r,+\infty)$ 上依次为正、负、正。
故前、后两个区间递增，中间递减。
\item 当 $k=1$ 时，$f_1'=x^2/(1+x)\ge0$，
且除 $x=0$ 外严格正。因此函数在整个
$(-1,+\infty)$ 上严格递增。
\item 当 $k>1$ 时，$-1<r<0$，导数在
$(-1,r),(r,0),(0,+\infty)$ 上依次为正、负、正。
故前、后两个区间递增，中间递减。
\end{enumerate}

\noindent\textbf{核验边界：}数学推导独立复核；
新源码尚未通过 XeLaTeX 编译验收。
% END_SOURCE_ENTRY path=derivative/reviewed/gaokao-2010-beijing-science-q18-solution.tex

% SOURCE_ENTRY path=derivative/reviewed/gaokao-2011-anhui-derivative-common-solution.tex sha=f43eab781a3dac98b956b8723ae2a67e2aecd4f9
% ORIGINAL_PREAMBLE \documentclass[UTF8,a4paper,11pt]{ctexart}
% ORIGINAL_PREAMBLE \usepackage{amsmath,amssymb}
% ORIGINAL_PREAMBLE \usepackage[margin=2.6cm]{geometry}
% ORIGINAL_PREAMBLE \usepackage[hidelinks]{hyperref}
% ORIGINAL_PREAMBLE \title{2011年安徽高考文理科共同导数题}
% ORIGINAL_PREAMBLE \author{三千弱水数学题库}
% ORIGINAL_PREAMBLE \date{2026年10月24日}
\clearpage
\phantomsection
\addcontentsline{toc}{section}{2011年安徽高考文理科共同导数题}
\section*{2011年安徽高考文理科共同导数题}
\noindent\textbf{出处：}文科第18题 \url{https://github.com/DxAThing/Gaokao-Math-Problems-Compilation/blob/f1ae169b58a75b3db6e35c670b3f0d5c78718188/content/2011/anhui_liberal.tex#L447}

\noindent 理科第16题 \url{https://github.com/DxAThing/Gaokao-Math-Problems-Compilation/blob/f1ae169b58a75b3db6e35c670b3f0d5c78718188/content/2011/anhui_science.tex#L362}

\paragraph{定义域和导数}
对 $a>0$，考虑函数
\[
f(x)=\frac{e^x}{1+ax^2},\qquad x\in\mathbb R.
\]
分母对全体实数严格为正，故
\[
f'(x)=\frac{e^x\bigl[a(x-1)^2+1-a\bigr]}{(1+ax^2)^2}.
\]

\paragraph{当 $a=4/3$ 时}
此时
\[
a(x-1)^2+1-a
=\frac43\left(x-\frac12\right)\left(x-\frac32\right).
\]
导数符号依次为正、负、正，
故极大值点为 $x=1/2$，极小值点为 $x=3/2$。
对应的极值分别是
\[
f(1/2)=\frac{3\sqrt e}{4},\qquad
f(3/2)=\frac{e^{3/2}}4.
\]

\paragraph{整条实数轴上的单调性}
当 $0<a<1$ 时，$a(x-1)^2+1-a>0$；
当 $a=1$ 时，该式为 $(x-1)^2$，
只在一个点为零，函数仍严格递增。
当 $a>1$ 时，导数在 $x=0$ 处为正、在 $x=1$ 处为负，
故不可能全局单调。
于是所求为
\[
\boxed{0<a\le1}.
\]

\bigskip
\noindent\textbf{质检说明：}数学推导独立验证通过；
本源码尚未实际编译，不得标注为编译通过。
% END_SOURCE_ENTRY path=derivative/reviewed/gaokao-2011-anhui-derivative-common-solution.tex

% SOURCE_ENTRY path=derivative/reviewed/gaokao-2013-fujian-science-q17-solution.tex sha=b4ef9241e47e073ce97c7b746e5a617552ac626d
% ORIGINAL_PREAMBLE \documentclass[UTF8,a4paper,11pt]{ctexart}
% ORIGINAL_PREAMBLE \usepackage{amsmath,amssymb}
% ORIGINAL_PREAMBLE \usepackage[margin=2.5cm]{geometry}
% ORIGINAL_PREAMBLE \usepackage[hidelinks]{hyperref}
% ORIGINAL_PREAMBLE \title{2013 年福建高考理科第 17 题\\切线与含参函数极值}
% ORIGINAL_PREAMBLE \date{2026 年 10 月 24 日}
\clearpage
\phantomsection
\addcontentsline{toc}{section}{2013 年福建高考理科第 17 题 · 切线与含参函数极值}
\section*{2013 年福建高考理科第 17 题\\切线与含参函数极值}
\noindent\textbf{原题来源：}\url{https://github.com/DxAThing/Gaokao-Math-Problems-Compilation/blob/f1ae169b58a75b3db6e35c670b3f0d5c78718188/content/2013/fujian_science.tex#L356}

设 $f_a(x)=x-a\ln x$，$x>0$，$a\in\mathbb R$。
\[
 f_a'(x)=1-\frac{a}{x}=\frac{x-a}{x}.
\]
\paragraph{当 $a=2$ 时的切线}
在点 $x=1$ 处，
$f_2(1)=1$、$f_2'(1)=-1$，
所以切线为
\[
 \boxed{y-1=-(x-1)}\quad\text{即}\quad x+y-2=0.
\]

\paragraph{极值的充要分类}
由于定义域中 $x>0$，导数分母为正：
\begin{enumerate}
\item 当 $a\le0$ 时，$x-a>0$，函数全域严格递增，
不存在极值点。
\item 当 $a>0$ 时，在 $(0,a)$ 导数为负，在
$(a,+\infty)$ 导数为正，因而 $x=a$ 是唯一极小值点，
极小值为
\[
 \boxed{f_a(a)=a-a\ln a}.
\]
在此情形也不存在极大值。
\end{enumerate}
\noindent\textbf{核验边界：}数学推导独立复核；
新源码尚未实际编译。
% END_SOURCE_ENTRY path=derivative/reviewed/gaokao-2013-fujian-science-q17-solution.tex

% SOURCE_ENTRY path=derivative/reviewed/gaokao-2014-anhui-common-q-solution.tex sha=32b6f2f195cc3c3a248fe4fb2191ca7d94fcfa64
% ORIGINAL_PREAMBLE \documentclass[UTF8,a4paper,11pt]{ctexart}
% ORIGINAL_PREAMBLE \usepackage{amsmath,amssymb}
% ORIGINAL_PREAMBLE \usepackage[margin=2.5cm]{geometry}
% ORIGINAL_PREAMBLE \usepackage[hidelinks]{hyperref}
% ORIGINAL_PREAMBLE \title{2014 年安徽高考文理共同题\\三次函数与边界参数质检}
% ORIGINAL_PREAMBLE \date{2026 年 10 月 24 日}
\clearpage
\phantomsection
\addcontentsline{toc}{section}{2014 年安徽高考文理共同题 · 三次函数与边界参数质检}
\section*{2014 年安徽高考文理共同题\\三次函数与边界参数质检}
\noindent\textbf{出处：}文科第 20 题 \url{https://github.com/DxAThing/Gaokao-Math-Problems-Compilation/blob/f1ae169b58a75b3db6e35c670b3f0d5c78718188/content/2014/anhui_liberal.tex#L628}

理科第 18 题 \url{https://github.com/DxAThing/Gaokao-Math-Problems-Compilation/blob/f1ae169b58a75b3db6e35c670b3f0d5c78718188/content/2014/anhui_science.tex#L279}

设 $f_a(x)=1+(1+a)x-x^2-x^3$，其中 $a>0$。
\paragraph{全实数的单调性}
\[
 f_a'(x)=1+a-2x-3x^2=-3(x-\alpha)(x-\beta)
\]
其中
\[
 \alpha=\frac{-1-\sqrt{4+3a}}3<0,\quad
 \beta=\frac{-1+\sqrt{4+3a}}3>0.
\]
导数符号由左至右为负、正、负。
所以函数在 $(-\infty,\alpha)$ 和
$(\beta,+\infty)$ 严格递减，在 $(\alpha,\beta)$ 严格递增。

\paragraph{$[0,1]$ 上最大值的取点}
由 $\beta>0$ 及
$\beta\le1\Longleftrightarrow a\le4$，可得
\[
 x_{\max}=
 \begin{cases}
  \dfrac{-1+\sqrt{4+3a}}3,&0<a<4,\\[2mm]
  1,&a\ge4.
 \end{cases}
\]

\paragraph{$[0,1]$ 上最小值的取点}
函数在该区间内先增后减或一直增，
内点不可能取得小于两端点的函数值。
比较 $f_a(0)=1$、$f_a(1)=a$，得到
\[
 x_{\min}=
 \begin{cases}
  1,&0<a<1,\\
  0\ \text{或}\ 1,&a=1,\\
  0,&a>1.
 \end{cases}
\]
特别地，$a=1$ 两个端点均取得最小值。

\paragraph{上游质检发现}
文科 LaTeX 的解析结尾缺少上式的完整分类，
理科版本解析完整。本稿保留两份原始卷别出处并独立补正。

\noindent\textbf{核验边界：}数学证明已独立复核；
新 TeX 尚未实际编译。
% END_SOURCE_ENTRY path=derivative/reviewed/gaokao-2014-anhui-common-q-solution.tex

% SOURCE_ENTRY path=derivative/reviewed/gaokao-2014-chongqing-liberal-q19-solution.tex sha=21c23d6ee39d2df72a866d88f5817108cfadfe79
% ORIGINAL_PREAMBLE \documentclass[UTF8,a4paper,11pt]{ctexart}
% ORIGINAL_PREAMBLE \usepackage{amsmath,amssymb}
% ORIGINAL_PREAMBLE \usepackage[margin=2.5cm]{geometry}
% ORIGINAL_PREAMBLE \usepackage[hidelinks]{hyperref}
% ORIGINAL_PREAMBLE \title{2014 年重庆高考文科第 19 题\\垂直切线与函数极值}
% ORIGINAL_PREAMBLE \date{2026 年 10 月 24 日}
\clearpage
\phantomsection
\addcontentsline{toc}{section}{2014 年重庆高考文科第 19 题 · 垂直切线与函数极值}
\section*{2014 年重庆高考文科第 19 题\\垂直切线与函数极值}
\noindent\textbf{原题来源：}\url{https://github.com/DxAThing/Gaokao-Math-Problems-Compilation/blob/f1ae169b58a75b3db6e35c670b3f0d5c78718188/content/2014/chongqing_liberal.tex#L560}

已知
\[
 f_a(x)=\frac{x}{4}+\frac{a}{x}-\ln x-\frac32,
 \qquad x>0,\quad a\in\mathbb R,
\]
在 $x=1$ 处的切线垂直于 $y=x/2$。
由
\[
 f_a'(x)=\frac14-\frac{a}{x^2}-\frac1x
\]
及垂直条件 $f_a'(1)=-2$，得
\[
 \frac14-a-1=-2,\qquad \boxed{a=\frac54}.
\]
此时
\[
 f'(x)=\frac{x^2-4x-5}{4x^2}
 =\frac{(x-5)(x+1)}{4x^2}.
\]
在 $x>0$ 的定义域内，$x+1$ 和 $4x^2$ 均为正，
故函数在 $(0,5)$ 上严格递减，
在 $(5,+\infty)$ 上严格递增。
极小值点为 $x=5$，极小值为
\[
 \boxed{f(5)=\frac54+\frac14-\ln5-\frac32=-\ln5}.
\]
不存在极大值；$x=-1$ 并非合法驻点。

\noindent\textbf{核验边界：}数学推导独立复核；
新源码尚未实际编译。
% END_SOURCE_ENTRY path=derivative/reviewed/gaokao-2014-chongqing-liberal-q19-solution.tex

% SOURCE_ENTRY path=derivative/reviewed/gaokao-2019-national2-liberal-q21-solution.tex sha=b6ee1b57ae209a8731d2f0cfe0ca2dea0b506bd3
% ORIGINAL_PREAMBLE \documentclass[11pt,a4paper,UTF8]{ctexart}
% ORIGINAL_PREAMBLE \usepackage{amsmath,amssymb}
% ORIGINAL_PREAMBLE \usepackage[margin=2.5cm]{geometry}
% ORIGINAL_PREAMBLE \usepackage[hidelinks]{hyperref}
% ORIGINAL_PREAMBLE \title{2019年全国Ⅱ卷（文科）第21题：导数与倒数零点}
% ORIGINAL_PREAMBLE \author{三千弱水数学题库·独立复核}
% ORIGINAL_PREAMBLE \date{2026年10月24日}
\clearpage
\phantomsection
\addcontentsline{toc}{section}{2019年全国Ⅱ卷（文科）第21题：导数与倒数零点}
\section*{2019年全国Ⅱ卷（文科）第21题：导数与倒数零点}
\noindent\textbf{出处：}2019年全国Ⅱ卷（文科）第21题。
参见上游固定版本：
\url{https://github.com/DxAThing/Gaokao-Math-Problems-Compilation/blob/f1ae169b58a75b3db6e35c670b3f0d5c78718188/content/2019/national_paper_2_liberal.tex#L418}。

\medskip
考察函数
\[
  f(x)=(x-1)\ln x-x-1,\qquad x>0.
\]

\paragraph{一、唯一极值点。}
定义域为 $(0,+\infty)$。求导得到
\[
 f'(x)=\ln x-\frac1x,\qquad
 f''(x)=\frac{x+1}{x^2}>0 \quad(x>0).
\]
因此 $f'$ 在整个定义域上严格递增。又
$f'(1)=-1<0$，$f'(e)=1-e^{-1}>0$。
故存在唯一 $\alpha\in(1,e)$ 使 $f'(\alpha)=0$。
于是 $f$ 在 $(0,\alpha)$ 上严格递减，在
$(\alpha,+\infty)$ 上严格递增；$\alpha$ 是唯一极值点，
且是严格极小值点。

\paragraph{二、恰有两个实根。}
注意 $f(1)=-2<0$，以及
\[
 \lim_{x\to0^+}f(x)=+\infty,
 \qquad \lim_{x\to+\infty}f(x)=+\infty.
\]
结合连续性及前述严格单调性，在 $(0,1)$ 内有且仅有一个根
$u$，在 $(\alpha,+\infty)$ 内有且仅有一个根 $v$；
其余位置不存在零点。

\paragraph{三、两根互为倒数。}
对每个 $x>0$，直接计算
\[
 f\!\left(\frac1x\right)
 =\left(\frac1x-1\right)\ln\frac1x-\frac1x-1
 =\frac{(x-1)\ln x-x-1}{x}
 =\frac{f(x)}{x}.
\]
故若 $f(u)=0$，则 $f(1/u)=0$。由 $0<u<1$，
$1/u>1$，而在 $(1,\alpha]$ 上 $f<0$，
故 $1/u=v$，因此
\[
 \boxed{uv=1}.
\]

\bigskip
\noindent\textbf{核验状态：}数学证明已独立复核。
本 TeX 源码尚未进行远端编译验收，不能标记为 \texttt{latex-compiled}。
原始试题内容及图片的版权请遵守上游权利声明。
% END_SOURCE_ENTRY path=derivative/reviewed/gaokao-2019-national2-liberal-q21-solution.tex

% SOURCE_ENTRY path=derivative/reviewed/gaokao-2023-beijing-q20-solution.tex sha=a4f54f7fbe2e6145b2a1075edf1c9d8890961420
% ORIGINAL_PREAMBLE \documentclass[UTF8,a4paper,11pt]{ctexart}
% ORIGINAL_PREAMBLE \usepackage{amsmath,amssymb}
% ORIGINAL_PREAMBLE \usepackage[margin=2.5cm]{geometry}
% ORIGINAL_PREAMBLE \usepackage[hidelinks]{hyperref}
% ORIGINAL_PREAMBLE \title{2023 年北京卷第 20 题：独立数学证明}
% ORIGINAL_PREAMBLE \date{2026 年 10 月 25 日}
\clearpage
\phantomsection
\addcontentsline{toc}{section}{2023 年北京卷第 20 题：独立数学证明}
\section*{2023 年北京卷第 20 题：独立数学证明}
\noindent\textbf{固定题源：}\url{https://github.com/DxAThing/Gaokao-Math-Problems-Compilation/blob/f1ae169b58a75b3db6e35c670b3f0d5c78718188/content/2023/beijing.tex#L627}

设 $f(x)=x-x^3e^{ax+b}$ 在 $x=1$ 处的切线为 $y=-x+1$。
\paragraph{第一问}
由 $f(1)=0,\ f'(1)=-1$ 和
$f'(x)=1-(3x^2+ax^3)e^{ax+b}$，
得 $e^{a+b}=1$、$(3+a)e^{a+b}=2$，故
\[\boxed{a=-1,\qquad b=1}.\]
\paragraph{第二问}
令 $g=f'$，则
\[
g(x)=1-x^2(3-x)e^{1-x},\qquad
g'(x)=-x(x^2-6x+6)e^{1-x}.
\]
记 $\alpha=3-\sqrt3$，$\beta=3+\sqrt3$。
因为指数因子恒正，$g'$ 在
$(-\infty,0)$、$(0,\alpha)$、$(\alpha,\beta)$、
$(\beta,+\infty)$ 上的符号依次为正、负、正、负。
因此 $g$ 在第一、三段严格递增，在第二、四段严格递减。

\paragraph{第三问}
有 $g(-1)=1-4e^2<0$，$g(0)=1>0$，
$g(1)=-1<0$，$g(3)=1>0$。
因 $1<\alpha$，又 $g$ 在 $(0,\alpha)$ 严格递减，
故 $g(\alpha)<g(1)<0$。结合连续性、严格单调性，
$g$ 在 $(-1,0)$、$(0,1)$、$(\alpha,3)$
上各有且仅有一个零点，其符号分别由负变正、由正变负、
由负变正；各单调区间内不可能出现额外零点。
且当 $x\ge3$ 时，$x^2(3-x)\le0$，从而 $g(x)\ge1>0$。
综上全部实数区间已覆盖，函数 $f$ 恰有
两个极小值点、一个极大值点：
\[
\boxed{\text{极值点共 }3\text{ 个}}.
\]

\noindent\textbf{审核状态：}数学推导已独立核验；
未进行 XeLaTeX 编译或 PDF 视觉质检。
% END_SOURCE_ENTRY path=derivative/reviewed/gaokao-2023-beijing-q20-solution.tex

% SOURCE_ENTRY path=derivative/reviewed/gaokao-2023-national-b-liberal-q20-solution.tex sha=905ae2ec7be7a86d9dda44ccf298845f73d5f154
% ORIGINAL_PREAMBLE \documentclass[UTF8,a4paper,11pt]{ctexart}
% ORIGINAL_PREAMBLE \usepackage{amsmath,amssymb}
% ORIGINAL_PREAMBLE \usepackage[margin=2.5cm]{geometry}
% ORIGINAL_PREAMBLE \usepackage[hidelinks]{hyperref}
% ORIGINAL_PREAMBLE \title{2023年全国乙卷文科第20题：切线与参数单调性}
% ORIGINAL_PREAMBLE \date{2026年10月24日}
\clearpage
\phantomsection
\addcontentsline{toc}{section}{2023年全国乙卷文科第20题：切线与参数单调性}
\section*{2023年全国乙卷文科第20题：切线与参数单调性}
\noindent\textbf{固定来源：}\url{https://github.com/DxAThing/Gaokao-Math-Problems-Compilation/blob/f1ae169b58a75b3db6e35c670b3f0d5c78718188/content/2023/national_paper_b_liberal.tex#L246}

设 $x>0$，
\[
f_a(x)=\left(\frac1x+a\right)\ln(1+x).
\]

\paragraph{一、切线}
\[
f_a'(x)=-\frac{\ln(1+x)}{x^2}
          +\frac{1/x+a}{1+x}.
\]
当 $a=-1$ 时，$f_{-1}(1)=0$，$f_{-1}'(1)=-\ln2$，
故曲线在 $x=1$ 处的切线为
\[
\boxed{y=-(\ln2)(x-1)}.
\]

\paragraph{二、单调递增的必要充分条件}
令
\[
G(x)=\frac{(1+x)\ln(1+x)-x}{x^2},\qquad x>0.
\]
则 $f_a'(x)=[a-G(x)]/(1+x)$。再令
\[
H(x)=\frac{x^2}{2}+x-(1+x)\ln(1+x).
\]
由 $H(0)=0$ 与 $H'(x)=x-\ln(1+x)>0$，得
$H(x)>0$，即
\[
G(x)<\frac12\quad(x>0).
\]
另一方面，用两次洛必达法则得到
\[
\lim_{x\to0^+}G(x)
 =\lim_{x\to0^+}\frac{\ln(1+x)}{2x}
 =\frac12.
\]
故 $\sup_{x>0}G(x)=1/2$，且上确界不取到。
当 $a\ge1/2$ 时，$f_a'(x)>0$ 恒成立；
当 $a<1/2$ 时，取充分小的 $x>0$ 有
$G(x)>a$，于是导数为负，不符合题意。
所以
\[
\boxed{a\in\left[\frac12,+\infty\right)}.
\]

\bigskip
\noindent\textbf{审校状态：}仅独立数学核验；新源码尚未实际
XeLaTeX 编译。原题及第三方图示版权归上游相应权利主体。
% END_SOURCE_ENTRY path=derivative/reviewed/gaokao-2023-national-b-liberal-q20-solution.tex

% SOURCE_ENTRY path=derivative/reviewed/gaokao-2023-new-gaokao-1-q19-solution.tex sha=43d5dab296813acd071e7ef109b15c850355f01b
% ORIGINAL_PREAMBLE \documentclass[UTF8,a4paper,11pt]{ctexart}
% ORIGINAL_PREAMBLE \usepackage{amsmath,amssymb}
% ORIGINAL_PREAMBLE \usepackage[margin=2.5cm]{geometry}
% ORIGINAL_PREAMBLE \usepackage[hidelinks]{hyperref}
% ORIGINAL_PREAMBLE \title{2023年新高考Ⅰ卷第19题：指数函数与严格不等式}
% ORIGINAL_PREAMBLE \date{2026年10月24日}
\clearpage
\phantomsection
\addcontentsline{toc}{section}{2023年新高考Ⅰ卷第19题：指数函数与严格不等式}
\section*{2023年新高考Ⅰ卷第19题：指数函数与严格不等式}
\noindent\textbf{固定来源：}\url{https://github.com/DxAThing/Gaokao-Math-Problems-Compilation/blob/f1ae169b58a75b3db6e35c670b3f0d5c78718188/content/2023/new_gaokao_paper_1.tex#L554}

设
\[
f_a(x)=a(e^x+a)-x,\qquad x\in\mathbb R,\ a\in\mathbb R.
\]
\paragraph{一、单调性}
$f_a'(x)=ae^x-1$。
当 $a\le0$ 时，导数在 $\mathbb R$ 上严格负，
$f_a$ 严格递减。
当 $a>0$ 时，唯一驻点
$x_*=-\ln a$，导数在其左侧为负、右侧为正。
因此 $f_a$ 先严格递减后严格递增，且
\[
\min_{x\in\mathbb R}f_a(x)=f_a(-\ln a)
=1+a^2+\ln a.
\]

\paragraph{二、证明全体 $x$ 上的严格不等式}
当 $a>0$ 时，令
\[
h(a)=a^2-\ln a-\frac12,\qquad
h'(a)=\frac{2a^2-1}{a}.
\]
由于导数由负变正，$h$ 在 $a=1/\sqrt2$ 处
取得唯一全局最小值，且
\[
h(a)\ge h(1/\sqrt2)=\ln\sqrt2>0.
\]
所以对所有实数 $x$ 均有
\[
\boxed{f_a(x)\ge1+a^2+\ln a
>2\ln a+\frac32}.
\]
不等式末端严格成立，包括函数自身取得极小值时。

\bigskip
\noindent\textbf{审校状态：}数学独立复核；
新 TeX 尚未编译。
% END_SOURCE_ENTRY path=derivative/reviewed/gaokao-2023-new-gaokao-1-q19-solution.tex

% SOURCE_ENTRY path=derivative/reviewed/gaokao-2024-two-proofs-solution.tex sha=3665f23093d454a1b74fd0b46dc7fbbcb5efee4e
% ORIGINAL_PREAMBLE \documentclass[UTF8,a4paper,11pt]{ctexart}
% ORIGINAL_PREAMBLE \usepackage[margin=25mm,headheight=15pt]{geometry}
% ORIGINAL_PREAMBLE \usepackage{amsmath,amssymb,amsthm,fancyhdr,hyperref,xcolor}
% ORIGINAL_PREAMBLE \hypersetup{colorlinks=true,urlcolor=blue!55!black,linkcolor=black}
% ORIGINAL_PREAMBLE \setlength{\parindent}{2em}
% ORIGINAL_PREAMBLE \setlength{\parskip}{0.25em}
% ORIGINAL_PREAMBLE \linespread{1.13}
% ORIGINAL_PREAMBLE \pagestyle{fancy}
% ORIGINAL_PREAMBLE \fancyhf{}
% ORIGINAL_PREAMBLE \fancyhead[L]{\small 三千弱水 · 高中导数真题}
% ORIGINAL_PREAMBLE \fancyhead[R]{\small 2024 年 · 首轮独立复核}
% ORIGINAL_PREAMBLE \fancyfoot[C]{\thepage}
% ORIGINAL_PREAMBLE \renewcommand{\headrulewidth}{0.3pt}
% ORIGINAL_PREAMBLE \newcommand{\qtitle}[1]{\par\bigskip\noindent{\large\bfseries #1}\par\medskip}
% ORIGINAL_PREAMBLE \newcommand{\sourceurl}[1]{\par\noindent{\small\color{black!65}来源：\href{#1}{固定版本 LaTeX 源码（点击查看）}}\par}
\clearpage
\phantomsection
\addcontentsline{toc}{section}{2024 年新高考导数两题：独立证明}
\section*{2024 年新高考导数两题：独立证明}
\begin{center}
  {\LARGE\bfseries 高考导数解答题：独立核验选编}\par\vspace{0.6em}
  {\large 2024 年新高考两道解答题}\par\vspace{0.8em}
  {\small 仅收录数学表达和独立推导；并非原试卷复刻}\par
\end{center}
\vspace{0.5em}
\noindent\textbf{审核说明：}已对下列两题给出的结论、关键代数推导、参数端点和定义域进行独立核查；未声称原试卷全文、上游仓库其他题目均正确。原题以所列源码链接为准。\par
\qtitle{一、2024 年新高考一卷第 18 题}
设
\[
 f(x)=\ln\frac{x}{2-x}+ax+b(x-1)^3,\qquad 0<x<2,
\]
其中 $a,b\in\mathbb R$。分别研究：(1) $b=0$ 且导数处处非负时 $a$ 的最小值；(2) 曲线的中心对称性；(3) $f(x)>-2$ 恰在 $1<x<2$ 成立时 $b$ 的范围。

\noindent\textbf{解答与核验。}

\noindent\textbf{(1)} 当 $b=0$ 时，
\[
 f'(x)=\frac{1}{x}+\frac{1}{2-x}+a=\frac{2}{x(2-x)}+a.
\]
由于 $x(2-x)=1-(x-1)^2\le 1$ 且 $x(2-x)>0$，有 $\frac{2}{x(2-x)}\ge 2$，等号在 $x=1$ 取到。因此 $f'(x)\ge0$ 对所有 $x\in(0,2)$ 成立当且仅当 $a\ge-2$，故 $a_{\min}=-2$。

\noindent\textbf{(2)} 注意到定义域 $(0,2)$ 在变换 $x\mapsto2-x$ 下不变，且
\[
 f(2-x)=-\ln\frac{x}{2-x}+a(2-x)-b(x-1)^3=2a-f(x).
\]
故 $\bigl(x,f(x)\bigr)$ 与 $\bigl(2-x,2a-f(x)\bigr)$ 均在曲线上，二者的中点为 $(1,a)$，曲线关于 $(1,a)$ 中心对称。

\noindent\textbf{(3)} 函数在 $x=1$ 连续，且题设要求 $x\le1$ 时 $f(x)\le-2$、$x>1$ 时 $f(x)>-2$。左右逼近 $1$ 可得 $f(1)=-2$，即 $a=-2$。记 $t=x-1\in(-1,1)$，则
\[
 f(1+t)+2=H_b(t):=\ln\frac{1+t}{1-t}-2t+bt^3.
\]
函数 $H_b$ 为奇函数，因而题设等价于 $H_b(t)>0$ 对一切 $0<t<1$ 成立。令
\[
 R(t)=\ln\frac{1+t}{1-t}-2t-\frac23t^3.
\]
对 $0<t<1$，
\[
 R'(t)=\frac{2}{1-t^2}-2-2t^2
       =\frac{2t^4}{1-t^2}>0,
 \qquad R(0)=0.
\]
于是 $R(t)>0$，且
\[
 H_b(t)=R(t)+\left(b+\frac23\right)t^3.
\]
若 $b\ge-\frac23$，则 $H_b(t)>0$。反之，
\[
 \lim_{t\to0^+}\frac{H_b(t)}{t^3}=b+\frac23,
\]
若 $b<-\frac23$，则充分小的正 $t$ 使 $H_b(t)<0$，矛盾。故 $\boxed{b\ge-\frac23}$。
\sourceurl{https://github.com/DxAThing/Gaokao-Math-Problems-Compilation/blob/f1ae169b58a75b3db6e35c670b3f0d5c78718188/content/2024/new_gaokao_paper_1.tex}

\qtitle{二、2024 年新高考二卷第 16 题}
设 $f(x)=e^x-ax-a^3$。分别求：(1) $a=1$ 时曲线在横坐标 $1$ 处的切线；(2) 存在严格负的极小值时参数 $a$ 的范围。

\noindent\textbf{解答与核验。}

\noindent\textbf{(1)} $a=1$ 时，$f'(x)=e^x-1$，$f(1)=e-2$，$f'(1)=e-1$。故切线为
\[
 y-(e-2)=(e-1)(x-1),\qquad\boxed{y=(e-1)x-1}.
\]
\noindent\textbf{(2)} 由 $f'(x)=e^x-a$、$f''(x)=e^x>0$，当 $a\le0$ 时 $f'(x)>0$，不存在极小值；当 $a>0$ 时唯一极小值点为 $x=\ln a$，极小值为
\[
 f(\ln a)=a-a\ln a-a^3=a\bigl(1-\ln a-a^2\bigr).
\]
因 $a>0$，极小值严格小于 $0$ 当且仅当 $G(a)=a^2+\ln a-1>0$。又
\[
 G'(a)=2a+\frac1a>0\quad(a>0),\qquad G(1)=0,
\]
所以所求为 $\boxed{a\in(1,+\infty)}$。
\sourceurl{https://github.com/DxAThing/Gaokao-Math-Problems-Compilation/blob/f1ae169b58a75b3db6e35c670b3f0d5c78718188/content/2024/new_gaokao_paper_2.tex}

\bigskip\noindent{\small\color{black!65}资料溯源：上游项目 DxAThing/Gaokao-Math-Problems-Compilation；其原创编辑与编排部分适用 CC BY-SA 4.0，部分原试题等第三方材料可能另有权利。本文推导为重新独立撰写，原项目与作者署名应继续保留。}
% END_SOURCE_ENTRY path=derivative/reviewed/gaokao-2024-two-proofs-solution.tex

% SOURCE_ENTRY path=derivative/reviewed/gaokao-2025-national2-q18-solution.tex sha=7c7c45a71c71e86597bb88ed17872bafbf1b5b0c
% ORIGINAL_PREAMBLE \documentclass[UTF8,a4paper,11pt]{ctexart}
% ORIGINAL_PREAMBLE \usepackage{amsmath,amssymb}
% ORIGINAL_PREAMBLE \usepackage[margin=2.5cm]{geometry}
% ORIGINAL_PREAMBLE \usepackage[hidelinks]{hyperref}
% ORIGINAL_PREAMBLE \title{2025年全国Ⅱ卷第18题：唯一极值、零点与对称差}
% ORIGINAL_PREAMBLE \date{2026年10月24日}
\clearpage
\phantomsection
\addcontentsline{toc}{section}{2025年全国Ⅱ卷第18题：唯一极值、零点与对称差}
\section*{2025年全国Ⅱ卷第18题：唯一极值、零点与对称差}
\noindent\textbf{固定来源：}\url{https://github.com/DxAThing/Gaokao-Math-Problems-Compilation/blob/f1ae169b58a75b3db6e35c670b3f0d5c78718188/content/2025/national_paper_2.tex#L496}

设 $0<k<1/3$，
\[
f(x)=\ln(1+x)-x+\frac{x^2}{2}-kx^3,\quad x>-1.
\]
仅在 $x>0$ 上讨论极值点 $x_1$ 和正零点 $x_2$。
\paragraph{一、唯一性}
\[
f'(x)=x^2\left(\frac1{1+x}-3k\right).
\]
令
\[
c=\frac1{3k}-1>0.
\]
对 $0<x<c$，$f'(x)>0$；对 $x>c$，$f'(x)<0$。
所以 $x_1=c$ 为唯一极值点，且是严格极大值点。
又 $f(0)=0$、$f(c)>0$，且
$\lim_{x\to\infty}f(x)=-\infty$。
由连续性及递减性，在 $(c,+\infty)$ 上有且仅有一个零点
$x_2$。注意 $x=0$ 不在题目要求的 $(0,+\infty)$ 内。

\paragraph{二、对称差严格递减}
令
\[
g(t)=f(c+t)-f(c-t),\qquad 0<t<c.
\]
链式法则给出正确的\textbf{加号}：
\[
g'(t)=f'(c+t)+f'(c-t).
\]
利用 $3k=1/(1+c)$，得
\[
f'(x)=\frac{x^2(c-x)}{(1+x)(1+c)}.
\]
令 $H(x)=x^2/(1+x)$，则
\[
H'(x)=\frac{x(x+2)}{(1+x)^2}>0\quad(x>0).
\]
因而
\[
g'(t)=\frac{t}{1+c}\bigl[H(c-t)-H(c+t)\bigr]<0,
\]
故 $g$ 在 $(0,c)$ 严格递减。

\paragraph{三、比较两零点位置}
由连续延拓的 $g(0)=0$ 及严格递减性，
\[
g(c)=f(2c)-f(0)=f(2c)<0.
\]
结合 $f(c)>0$，及 $f$ 在 $(c,\infty)$ 严格递减，
唯一正零点 $x_2$ 满足
\[
\boxed{x_1<x_2<2x_1}.
\]

\bigskip
\noindent\textbf{源码质检提示：}
上游解答中一行把 $g'(t)$ 误写为两个 $f'$ 之差；
链式法则正确结果是上文的\textbf{和}。
本稿为独立数学证明，尚未实际 XeLaTeX 编译。
% END_SOURCE_ENTRY path=derivative/reviewed/gaokao-2025-national2-q18-solution.tex

% SOURCE_ENTRY path=derivative/reviewed/gaokao-2025-shanghai-q19-solution.tex sha=9bbe376e6b111189714f02ea2075ed9fb4f285c7
% ORIGINAL_PREAMBLE \documentclass[UTF8,a4paper,11pt]{ctexart}
% ORIGINAL_PREAMBLE \usepackage{amsmath,amssymb}
% ORIGINAL_PREAMBLE \usepackage[margin=2.5cm]{geometry}
% ORIGINAL_PREAMBLE \usepackage[hidelinks]{hyperref}
% ORIGINAL_PREAMBLE \title{2025年上海卷第19题：含参对数函数的极大值}
% ORIGINAL_PREAMBLE \date{2026年10月24日}
\clearpage
\phantomsection
\addcontentsline{toc}{section}{2025年上海卷第19题：含参对数函数的极大值}
\section*{2025年上海卷第19题：含参对数函数的极大值}
\noindent\textbf{固定来源：}\url{https://github.com/DxAThing/Gaokao-Math-Problems-Compilation/blob/f1ae169b58a75b3db6e35c670b3f0d5c78718188/content/2025/shanghai.tex#L381}

设
\[
 f_m(x)=x^2-(m+2)x+m\ln x,\quad x>0,\ m\in\mathbb R.
\]
\paragraph{一、求不等式解集}
$f_m(1)=-m-1=0$ 导致 $m=-1$。
此时 $f_{-1}(x)\le x^2-1$ 等价于
$x+\ln x\ge1$。因为
\[
\frac{\mathrm d}{\mathrm dx}(x+\ln x)=1+\frac1x>0
\quad(x>0)
\]
且 $1+\ln1=1$，故解集为
\[
\boxed{[1,+\infty)}.
\]

\paragraph{二、极大值存在的参数}
求导并分解为
\[
f_m'(x)=\frac{(2x-m)(x-1)}{x}.
\]
\begin{enumerate}
\item $m\le0$ 时，合法驻点仅 $x=1$，
导数由负变正，仅有极小值。
\item $0<m<2$ 时，$0<m/2<1$，
导数正、负、正，在 $x=m/2$ 取得严格极大值。
\item $m=2$ 时，$f_m'(x)=2(x-1)^2/x\ge0$，
函数严格递增，没有极大值。
\item $m>2$ 时，$1<m/2$，导数正、负、正，
在 $x=1$ 取得严格极大值。
\end{enumerate}
故参数范围为
\[
\boxed{m\in(0,2)\cup(2,+\infty)}.
\]
特别注意 $m=2$ 的重驻点不是极值点。

\bigskip
\noindent\textbf{审校状态：}数学独立复核；
新源码尚未实际编译。
% END_SOURCE_ENTRY path=derivative/reviewed/gaokao-2025-shanghai-q19-solution.tex

\end{document}
